% ---
% titre: Problème de proportionnalité - la vente de pâtisserie
% theme: FA
% chapitre: Proportionnalité
% sous_chapitre: Résoudre des problèmes
% niveau: [10e]
% type: EVAL
% tags: [proportionnalite,c-nombres-decimaux, p-question-TS]
% difficulte: 1
% corrige: true
% ---

\question[2]
Au marché, 5 barquettes de fraises coûtent 5,75 francs.

\begin{parts}
\vspace{20pt}\part[2] Calcule le prix de 4 barquettes de fraises.

\vspace{10pt}{\footnotesize\raggedleft Espace pour ta démarche et tes calculs\par}
\begin{solutionorgrid}[140pt]

%\vspace{5pt}
%\def\arraystretch{1.5}
%\begin{tabularx}{0.5\textwidth}{l|C|C|}
%Somme totale [CHF]&830&\\
%\hline
%Somme par élève [CHF]&41,5&50\\
%\end{tabularx}
%\def\arraystretch{1.5}
%
%\hfill\textbf{(0.5pt pour la construction du tableau)}
%
%\vspace{15pt}$50 \cdot 830 = 41\,500$ \hfill\textbf{(0.5pt)}
%
%\vspace{5pt}$41\,500 : 41,5 = 1000$\hfill\textbf{(0.5pt)}

\end{solutionorgrid}

\begin{tcolorbox}[enhanced jigsaw, interior hidden, colframe=box, parbox=false]%
Réponse: 
\sol{\vspace{20pt}}{Il aurait fallu gagner 1\,000.- pour que chaque élève reçoive 50.- \textbf{(0.5pt)}}
\end{tcolorbox}

\vspace{20pt}\part[2] Combien peut-on acheter de barquettes avec 12,65 francs?

\vspace{10pt}{\footnotesize\raggedleft Espace pour ta démarche et tes calculs\par}
\begin{solutionorgrid}[140pt]
\end{solutionorgrid}

\begin{tcolorbox}[enhanced jigsaw, interior hidden, colframe=box, parbox=false]%
Réponse: 
\sol{\vspace{20pt}}{Il aurait fallu gagner 1\,000.- pour que chaque élève reçoive 50.- \textbf{(0.5pt)}}
\end{tcolorbox}



\end{parts}


